\section{A global Cayley subgroup-curve energy bound}
\label{sec:cayley-curve-uniform}
The scaling-rigidity result above explains why the two-parameter collision
family cannot be reduced to one polynomial by elementary coordinate changes.
There is nevertheless a uniform bound for every subgroup size, obtained by
applying the positive-characteristic curve theorem of Corvaja--Zannier~\cite{CorvajaZannier2013},
\emph{J. Eur. Math. Soc.} \textbf{15} (2013), Theorem~2/Corollary~2,
to each off-diagonal ratio curve.  The exponent remains
$8/3$, but the earlier parity and size restrictions disappear.

Let $p$ be an odd prime, let $H\leq\mathbf F_p^*$ have order $K\geq2$, and put
\[
 c(h)=\frac{h+1}{h-1},\qquad A=c(H\setminus\{1\}),
\]
with ratio-fibre map
\[
 d_\rho(x)=\frac{(\rho x-1)(x-1)}{(1-\rho)x},
 \qquad \rho\in H\setminus\{1\}.
\]
For $\rho,\sigma\in H\setminus\{1\}$ define the extended collision count
\[
 N(\rho,\sigma)=\#\{(x,y)\in H^2:d_\rho(x)=d_\sigma(y)\}.
\]
The original pair-key energy satisfies
\begin{equation}\label{eq:czu-ec}
 E_C:=\sum_t C_H(t)^2=E_+(A)-(K-1)^2.
\end{equation}

\begin{theorem}[Uniform Cayley curve bound]\label{thm:czu-cayley}
For every odd prime $p$, every subgroup $H\leq\mathbf F_p^*$ of order
$K\geq2$, and every distinct $\rho,\sigma\in H\setminus\{1\}$,
\begin{equation}\label{eq:czu-point}
 N(\rho,\sigma)\leq
 \min\!\left\{2K,
 \max\!\left(3\sqrt[3]{32}\,K^{2/3},\frac{48K^2}{p}\right)\right\}.
\end{equation}
If
\[
 B(K,p)=\min\!\left\{2K,
 \max\!\left(3\sqrt[3]{32}\,K^{2/3},\frac{48K^2}{p}\right)\right\}
\]
and
\[
 D_K=2K^2-7K+4+\gcd(2,K),
\]
then
\begin{equation}\label{eq:czu-energy-exact}
 E_+(A)-(K-1)^2
 \leq D_K+(K-1)(K-2)\bigl(B(K,p)-4\bigr).
\end{equation}
Consequently, uniformly for every subgroup size,
\begin{equation}\label{eq:czu-energy-asymp}
 E_+(A)\ll K^{8/3}+\frac{K^4}{p},
\end{equation}
with an absolute implied constant.
\end{theorem}

\begin{proof}
Clearing denominators in $d_\rho(x)=d_\sigma(y)$ gives
\begin{equation}\label{eq:czu-F}
 F_{\rho,\sigma}(x,y)
 =(1-\sigma)y(\rho x-1)(x-1)
 -(1-\rho)x(\sigma y-1)(y-1).
\end{equation}
For $\rho\neq\sigma$ this has bidegree $(2,2)$.  Viewed as a quadratic in
$y$, its discriminant is a quartic in $x$, and direct elimination gives
\begin{equation}\label{eq:czu-disc}
 \operatorname{disc}_x\!\bigl(\operatorname{disc}_yF_{\rho,\sigma}\bigr)
 =4096\rho^2\sigma^2(\rho-1)^6(\rho-\sigma)^2(\sigma-1)^6.
\end{equation}
The right side is nonzero in odd characteristic for distinct nonzero
nonunit parameters.  Hence the quadratic discriminant is squarefree and is
not a square in $\overline{\mathbf F}_p(x)$; therefore the curve is
geometrically irreducible.

Its closure in $\mathbf P^1\times\mathbf P^1$ has exactly the four boundary
points
\[
 (0,0),\quad(0,\infty),\quad(\infty,0),\quad(\infty,\infty).
\]
In the corresponding affine charts the pairs of first derivatives are
\[
 (\rho-1,1-\sigma),\quad
 (\sigma(\rho-1),1-\sigma),\quad
 (\rho-1,\rho(1-\sigma)),\quad
 (\sigma(\rho-1),\rho(1-\sigma)),
\]
so none is singular.  Equivalently, \eqref{eq:czu-disc} exhibits the
normalization as a double cover of $\mathbf P^1$ branched at four points,
thus of genus one.  Since a bidegree $(2,2)$ divisor has arithmetic genus
one, there are no hidden projective singularities.

The coordinate functions $x$ and $y$ both have degree two.  Their valuation
pairs at the four boundary points are
\[
 (1,1),\quad(1,-1),\quad(-1,1),\quad(-1,-1),
\]
which proves multiplicative independence modulo constants.  With $S$ the
boundary divisor, the affine Euler characteristic entering
Corvaja--Zannier is therefore
\[
 \chi=|S|+2g-2=4.
\]
The degree-two coordinate maps are separable because $p$ is odd, while
$p\nmid K$ because $K\mid p-1$.  Corollary~2 of that paper~\cite{CorvajaZannier2013}, applied with
$m_1=m_2=K$, $d_1=d_2=2$, and $\chi=4$, now yields
\[
 N(\rho,\sigma)
 \leq\max\!\left(3\sqrt[3]{32}\,K^{2/3},\frac{48K^2}{p}\right).
\]
For each fixed $x\in H$, equation \eqref{eq:czu-F} is quadratic in $y$,
which also gives the trivial bound $N(\rho,\sigma)\leq2K$.  This proves
\eqref{eq:czu-point}.

It remains to restore the exact pair-key domain.  Extending one ratio fibre
to all $x\in H$ adds precisely $x=1$ and $x=\rho^{-1}$, and both have value
zero.  Hence every ordered parameter pair contributes exactly four
artificial zero--zero collisions, giving the exact identity
\begin{equation}\label{eq:czu-restore}
 E_C=\sum_{\rho,\sigma\neq1}N(\rho,\sigma)-4(K-1)^2.
\end{equation}
On the diagonal,
\[
 F_{\rho,\rho}(x,y)=-(\rho-1)(x-y)(\rho xy-1).
\]
Summing the intersections of these two components over $\rho\neq1$ gives
$K-\gcd(2,K)$ and therefore
\[
 \sum_{\rho\neq1}N(\rho,\rho)
 =2K(K-1)-K+\gcd(2,K).
\]
After subtracting the four artificial collisions for every diagonal ratio,
the total diagonal contribution to $E_C$ is exactly $D_K$.  There are
$(K-1)(K-2)$ ordered off-diagonal parameter pairs; inserting
\eqref{eq:czu-point} into \eqref{eq:czu-restore} proves
\eqref{eq:czu-energy-exact}.  Finally
$B(K,p)\ll K^{2/3}+K^2/p$, and the diagonal term is $O(K^2)$, which gives
\eqref{eq:czu-energy-asymp}.
\end{proof}

\begin{remark}[Range and sharp frontier]
The characteristic-free term dominates whenever
$128K^4\leq p^3$, equivalently
$K\leq2^{-7/4}p^{3/4}$.  Its coefficient is
$3\sqrt[3]{32}=9.5244063118\ldots$.  This removes the old odd-order and
subgroup-size restrictions from the project curve route, but it does not
improve the exponent.  The centered pointwise target
\[
 E_+(A)=\frac{(K-1)^4}{p}+O\!\left(K^{2+o(1)}\right)
\]
remains open and requires cancellation across the two-parameter family,
not merely a worst-curve estimate.
\end{remark}

\begin{remark}[Exact regression boundary]
The maintained exact verifier checks the discriminant identity, all four
boundary gradients, diagonal factorization, restoration correction and the
integer form of the energy estimate for every subgroup order $K\geq3$
occurring for odd primes $p\leq257$.  This is a finite algebra and
implementation regression; it does not re-prove the external Corvaja--Zannier theorem.
\end{remark}
